This review examines how extensions of classical continuum mechanics change thermomechanical wave propagation, energy transport and interface phenomena. It organizes generalized thermoelastic, multiphysics, memory-dependent, fractional, gradient and nonlocal models by constitutive mechanism, governing structure, predicted modes and available validation. The literature map is a purposive inventory assembled from the repository’s existing collection and targeted bibliographic checks; DOI metadata checks are distinguished from full-text reading and claim-level verification. Several corrections follow from balancing conventions and explicit model equations. With total energy per unit volume and outward-positive flux, the classical mechanical contribution is minus stress power, so that $J_i=q_i-\sigma_{ij}v_j$. Lord–Shulman and Green–Lindsay models are not conflated with two-temperature theories; the former’s relaxation acts on the heat flux, whereas Green–Lindsay introduces temperature-rate dependence in a single-temperature formulation. A standard Eringen-type nonlocal dispersion relation has a bounded high-wavenumber frequency, not an unbounded frequency or group speed. Quasi-electrostatic piezoelectric models do not contain a propagating electromagnetic branch. Generalized-theory energy fluxes, benchmark values and several literature-dependent claims remain conditional or open where their derivations or primary-source checks are incomplete. The review therefore separates derived results, metadata checks, secondary-source findings and unresolved items, and does not present an unbuilt revision as a validated manuscript.